\subsection{\texorpdfstring{COMPLETE SUBSETS OF \(\mathcal{F}_{\mathfrak{S}}(\mathbf{X};\mathbf{Y})\)}{COMPLETE SUBSETS OF F SUB S (X;Y)}}

PROPOSITION 5. Let \(\Phi\) be a set, Y a uniform space and \(\mathfrak{S}\) a set of subsets of X. Then a filter \(\Phi\) on \(\mathcal{F}_{\mathfrak{S}}(\mathbf{X};\mathbf{Y})\) converges to \(u_0\) if and only if \(\Phi\) is a Cauchy filter with respect to the uniformity of \(\mathfrak{S}\) -convergence and converges pointwise to \(u_0\) in \(\mathbf{B} = \bigcup_{\mathbf{A}\in \mathfrak{S}}\mathbf{A}\) .

Since the structure of pointwise convergence in B is coarser than that of S-convergence, it is enough to show that for each \(A \in S\) and each closed entourage V of Y, \(W(A, V)\) is closed in B with respect to the topology of pointwise convergence (Chapter II, § 3, no. 3, Proposition 7). Now \(W(A, V)\) is the intersection of the inverse images of V under the mappings \((u, v) \to (u(x), v(x))\) as x runs through A; these mappings are continuous with respect to the topology of pointwise convergence (no. 2, Remark 6), and the result follows.

COROLLARY 1. A subspace H of \(\mathcal{F}_{\mathfrak{S}}(\mathbf{X};\mathbf{Y})\) is complete if and only if, for each Cauchy filter \(\Phi\) on H, there exists \(u_0\in \mathbf{H}\) such that \(\Phi\) converges pointwise to \(u_0\) in \(\mathbf{B} = \bigcup_{\mathbf{A}\in \mathfrak{S}}\mathbf{A}\) .

This follows immediately from Proposition 5.

COROLLARY 2. Let \(\mathfrak{S}_1, \mathfrak{S}_2\) be two sets of subsets of X, whose union is the same and which are such that \(\mathfrak{S}_1 \subset \mathfrak{S}_2\) , and let H be a subset of \(\mathcal{F}(X; Y)\) . Then if H is complete with respect to \(\mathfrak{S}_1\) -convergence, it is complete with respect to \(\mathfrak{S}_2\) -convergence.

For every Cauchy filter with respect to \(S_{2}\) -convergence is also a Cauchy filter with respect to \(S_{1}\) -convergence, and we may apply Corollary 1.

COROLLARY 3. Let H be a subset of \(\mathcal{F}(\mathbf{X};\mathbf{Y})\) such that, for each

\[
x \notin \mathrm{B} = \bigcup_ {\mathrm{A} \in \mathfrak {S}} \mathrm{A},
\]

the closure of \(\mathbf{H}(x)\) in \(\mathbf{Y}\) is a complete subspace of \(\mathbf{Y}\) . Then the closure \(\overline{\mathbf{H}}\) of \(\mathbf{H}\) in \(\mathcal{F}_{\mathfrak{S}}(\mathbf{X};\mathbf{Y})\) is a complete subspace.

Let \(\Phi\) be a Cauchy filter on \(\overline{\mathbf{H}}\) , and define a mapping \(v: \mathbf{X} \to \mathbf{Y}\) as follows. If \(x \in \mathbf{B}\) , \(\Phi(x)\) is a Cauchy filter on \(\overline{\mathbf{H}(x)}\) (no. 2, Remark 6), hence by hypothesis it has at least one limit point; take \(v(x)\) to be one of these limits. If \(x \notin \mathbf{B}\) , take \(v(x)\) to be any point of \(\mathbf{Y}\) . With this definition of \(v\) , it is clear that \(\Phi\) converges pointwise to \(v\) in \(\mathbf{B}\) , and \(v\) is therefore a limit of \(\Phi\) in \(\mathcal{F}_{\mathfrak{S}}(\mathbf{X}; \mathbf{Y})\) by Proposition 5.

In particular, if Y is complete, the hypothesis of Corollary 3 of Proposition 5 is satisfied for every \(H \subset \mathcal{F}(X; Y)\) ; hence:

THEOREM 1. Let X be a set, let S be a set of subsets of X, and let Y be a complete uniform space. Then the uniform space \(\mathcal{F}_{\mathfrak{S}}(X;Y)\) is complete.
% semantic-fragments: legacy-input-seam
\relax{}
